\section{Discussion}
\label{sec:discussion}

\subsection{The replay-point cliff is not an optimization answer}

The cliff measurement is deterministic: only post-pool replay is smaller than
the raw input for this architecture. The primary results change its
interpretation. Compression determines how many exemplars \emph{fit}, but not
how useful the remaining trainable subspace is. Maximum post-pool replay (A1)
fixes the representation and buys volume; raw replay is expensive per item but
leaves budget for attention-level updates. Across all 21 patients, the second
allocation improves more patients (16/21 versus 12/21) and every pooled
minority-class metric, even though it stores an order of magnitude fewer
exemplars. The cliff tells you where you \emph{may} replay; it does not tell you
where you \emph{should}.
\footnote{All primary-cohort figures use the pre-registered locked estimator
(median over seeds within patient, then across patients). Secondary and
sensitivity analyses (E9 ablation, E15 regularization, E16 reserve, E17
split-first, E12 label curve, E13 cross-database) report seeds-mean-averaged
descriptive statistics; those values are flagged where they appear and their
median unification is deferred (\texttt{BLOCKERS.md}, ESC-4).}

\subsection{Why encoder adaptation may help}

The source model has strong V sensitivity but weak S sensitivity. S
identification depends heavily on prematurity and RR context, yet the base
architecture fuses only a two-dimensional RR projection \emph{after} pooling. An
encoder LoRA update can reshape morphology--context interactions upstream of the
pool that a replacement linear head cannot reach. The F-class result is
consistent: post-pool arms plateau near $0.39$ F1 while encoder arms reach
$0.58$--$0.60$. This remains a mechanistic \emph{hypothesis}. RQ5 is the test that
speaks to it most directly, and it is weakly consistent with the hypothesis
rather than confirmatory: starting from sources whose S sensitivity is already
$0.78$--$0.88$ \emph{on their own DS1 test split}, encoder adaptation still
matches head-only replay ($\Delta_{\mathrm{A4-A1}} = +0.0001$ and $+0.0024$).
The adaptation gain persists under two alternative class-reweighted
source-training objectives. This supports robustness to the choice of source
objective, but does not establish robustness to a demonstrably stronger
unseen-patient source model. Those margins are also small and carry no
reported uncertainty, so they do not establish an ordering.

\subsection{What the data support}

Supported by the 21-record patient-disjoint, five-seed primary evidence:
\begin{itemize}[leftmargin=*]
\item patient-specific, target-informed adaptation improves the primary cohort
under the analytical 16\,KiB budget;
\item both encoder-LoRA arms significantly beat the frozen model at patient level
(Holm $p\le0.017$), with large rank-biserial effects and CIs excluding zero ---
these are the \emph{only} two of the six pre-specified comparisons that survive
Holm correction;
\item encoder arms improve more patients and give better pooled S and F
performance than maximum head-only replay, with $\approx 12\times$ fewer
exemplars, though the margin in improved-patient count narrows from four patients
to two once a meaningful-change threshold is applied;
\item the target-side benefit is attributable to encoder plasticity (B4$-$B2,
$p = 0.005$, 16/21 patients); the replay contrast is positive in mean but not
statistically supported (B4$-$B3, $p = 0.473$, median $\approx$ zero), and replay
\emph{volume} shows \textbf{no statistically supported target-side benefit
under the tested selector, replay ratio, optimizer, and patient cohort}
(B1$-$B6, 5 improved / 15 worsened);
\item replay provides substantially better source retention than the tested
zero-byte $B$-factor anchor under the evaluated configuration: at matched bytes,
replay gives $+0.136$ target gain against the anchor's $+0.078$ while losing
about a fifth as much source performance ($-0.027$ versus $-0.136$ \macrof{}),
and the anchor adds nothing once replay is present. Only this one
scale-ambiguous regularizer was tested, so this does not speak to EWC,
distillation, output consistency, or a $\|AB\|_F^2$ penalty
(Section~\ref{sec:reg});
\item about 100 labeled beats recover $\approx 80\%$ of the \emph{cohort-mean}
500-label gain, and whether the adaptation window covers the patient's future
arrhythmia class matters more than the label count (RQ7);
\item the conclusions do not depend on saturating the budget: re-running the
whole grid against a $15{,}360$\,B ceiling, i.e.\ holding back a 1\,KiB firmware
reserve, preserves every significant gain over the frozen model and keeps
encoder adaptation among the top allocations. The \emph{exact arm ordering is
not preserved} --- A1 rises to $0.793$ while A5 falls to $0.792$, putting A4
first --- which is why we read rank-1 as the more robust operating point
(Section~\ref{sec:reserve}).
\end{itemize}

Supported by secondary evidence, on reduced panels or other cohorts:
\begin{itemize}[leftmargin=*]
\item (RQ4) within the measured 12-record Adam panel at ranks 1 and 2, raising
the adaptation-state budget from 8 to 64\,KiB produced no statistically
detectable benefit; 14 of 18 fixed-configuration contrasts met the
pre-specified $\pm 0.05$ equivalence criterion, uncorrected for multiplicity;
\item the adaptation gain survives adaptation from two class-reweighted DS1
sources, of five trained; the encoder-minus-head margins there are nonnegative,
but they are small and carry no reported uncertainty (RQ5);
\item chronological target-informed adaptation transfers to two external
databases at \emph{subject} level, with $+0.14$ to $+0.24$ mean gains on INCART
($6/6$ subjects) and SVDB ($45$--$46/46$); encoder LoRA has the higher mean on
both external subsets, without a paired direct-comparison interval (RQ6).
\end{itemize}

Not supported:
\begin{itemize}[leftmargin=*]
\item direct statistical superiority of A4/A5 over A1 on the primary cohort (the
direct comparisons remain non-significant, and the reweighted-source and
external margins, while consistently positive, are small);
\item any optimizer ranking, or any rank-4/rank-8 statement: SGD is measured only
at 8\,KiB, SGD-with-momentum not at all, and 86 of 140 planned sweep cells are
unrun;
\item \emph{inferential} external validity on INCART: the subject-level gain is
positive with an interval excluding zero, but $n = 6$ subjects supports a
descriptive reading only, and the cohort is a tractability-capped sample of 13 of
75 records rather than a principled selection;
\item RR-architecture and CNN source baselines (three E11 variants still require
new model graphs);
\item any hardware runtime, SRAM, latency, or energy claim;
\item a ``50 labels are sufficient'' claim, or any per-patient reading of the
label curve: at 100 labels only $38\%$ of patients individually reach $80\%$
recovery, so the cohort-mean figure must not be read as a per-patient guarantee.
\end{itemize}

\subsection{Embedded interpretation}

Everything reported here is persistent-state accounting (Appendix~\ref{sec:embedded}
gives the packed record layouts, the memory-category table, and a qualitative
discussion of the unmatched compute paths): trainable weights,
gradients, optimizer moments, replay values, labels, and metadata, packed and
asserted per configuration. It excludes forward and backward activations, stack
and workspace, allocator overhead, and DMA buffers, all of which are peak-SRAM
rather than persistent-state quantities and all of which are implementation- and
toolchain-specific. The 16\,KiB figure is therefore a statement about what a
device must \emph{retain} between adaptation steps, not about what it must have
free at any instant; a hardware study would need to add the runtime terms.

The budget is also accompanied by a firmware implementation reserve. An arm
consuming 16{,}383 of 16{,}384 bytes is not shippable, so we re-generated all
configurations against an effective ceiling of 15{,}360\,B and re-ran every arm
at the reduced replay counts rather than reporting old numbers against new byte
figures. That costs A1 49 exemplars and the encoder arms five each. It changes no
significance conclusion, but it does reshuffle the arm ordering --- A5 is the
arm most sensitive to losing its last few replay samples --- so the reserve
replication is what identifies rank-1 as the safer engineering choice
(Section~\ref{sec:reserve}).

There is one cost of the encoder allocation that the byte budget hides. Raw
replay requires re-forwarding stored beats through the encoder and propagating
gradients through the LoRA training path, whereas post-pool replay enters the
graph after the encoder and bypasses it. Encoder adaptation therefore has
substantially greater compute and energy cost per adaptation step, but
\textbf{this study does not measure or reliably estimate its magnitude} and we
quote no ratio (Appendix~\ref{sec:embedded} explains why the earlier analytical
figure was withdrawn). The arms are matched in bytes and not in compute, which on
a duty-cycled device is an energy difference a deployment study would have to
quantify before choosing between them.

The external result adds a nuance to the frozen-replay story: under strong domain
shift, the no-replay encoder LoRA achieves the \emph{largest} target gain but the
\emph{most} source-domain change. Replay is therefore best understood not as a pure
accuracy lever but as a stability regulator whose value grows with the cost of
source-domain loss.

The regularization comparison sharpens that reading. A parameter-space anchor and
a replay buffer are both stability mechanisms, but they are not
interchangeable: the anchor constrains \emph{how far the weights may move},
whereas replay constrains \emph{what the function must still do}. On a shift that
is a change of patient rather than a change of task, the second constraint is the
useful one, and the measurements bear that out --- at equal persistent bytes,
replay incurs roughly one fifth of the anchor's source-domain loss while also
adapting further. We hold this to its tested scope: \emph{replay provides
substantially better source retention than the tested zero-byte $B$-factor
anchor under the evaluated configuration}. It is one regularizer, in a
scale-ambiguous form, so it does not license the broader claim that a budget
should never be spent on regularization --- restricted EWC, source-logit
distillation, output-consistency regularization, and a $\|AB\|_F^2$ penalty are
all untested here and could behave differently.

\subsection{Backprop memory is not negligible at the KiB scale}

A recent survey of on-device learning in TinyML states the prevailing
assumption directly: ``since all continual learning solutions maintain a replay
buffer of sufficient size to prevent catastrophic forgetting, the memory cost
of backpropagation is unlikely to constitute a bottleneck relative to the buffer
itself'' \cite{pavan2026odlsurvey}. Our two-axis account contradicts this in the
16\,KiB regime. The replay buffer is bounded at 16\,KiB by construction, but the
peak transient working set of the raw-replay encoder arms is $41.3\kib$ even
under block rematerialisation at FP16 (Section~\ref{sec:method},
Appendix~\ref{sec:embedded})---roughly $2.5\times$ the buffer, not a fraction of
it. The assumption holds when the frozen frontend is deep and the buffer is
large; it fails precisely at the KiB scale and pooled-transformer topology
studied here, where a single $8{,}448$-element feed-forward tensor and its
gradient dominate. The correct statement is conditional: backprop memory is
negligible against the buffer only when the trainable subgraph is shallow (as in
post-pool adaptation, whose transient is under $0.1\kib$); for encoder-level
adaptation it is the larger of the two transient charges, and whether it is
affordable depends on how much of it aliases the inference arena the device
already provisions.

There is a corollary for the standard memory-for-compute trade. Block-level
rematerialisation buys the encoder arms a $1.87\times$ peak-transient reduction
($79{,}164 \to 42{,}312\bytes$) at a $1.993\times$ forward-compute multiplier
(\texttt{results/e6\_two\_axis\_accounting.csv}), so at this scale you pay
marginally more in compute than you recover in memory and checkpointing is close
to value-neutral. The reason is structural: the $8{,}448$-element feed-forward
hidden tensor and its gradient dominate the working set so completely that
checkpointing at block granularity cannot escape them, which implies that further
shrinking the encoder-adaptation transient requires finer-grained
rematerialisation or a topology change, not more block checkpointing.

\subsection{Responder structure (exploratory, post-hoc)}

\emph{This subsection is exploratory and post-hoc. It was not pre-registered, it
is hypothesis-generating rather than confirmatory, and it would need its own
pre-registration to be treated as a test. No inferential statistic is reported for
it, and it does not change the pre-registered primary result.}

Looking at where the A4-versus-A1 contrast lives across patients is descriptively
informative. The per-patient median-within A4$-$A1 differences
(\texttt{results/r8\_responder\_exploratory.csv}) split 14 better, 3 exactly
tied, and 4 worse, and range from $-0.033$ to $+0.167$; three patients alone
account for $61\%$ of the total positive mass. In other words the two allocations
look interchangeable for most patients and materially apart for a small
minority. This single observation is consistent, descriptively, with the whole
cluster of otherwise puzzling numbers in Table~\ref{tab:pairwise_tests}: a median
paired difference near zero, a rank-biserial around $0.72$, a mean paired
difference an order of magnitude above the median, and an equivalence test that
fails on its upper bound alone. A natural next question---whether the responding
minority is characterised by higher minority-class support in the adaptation
window---cannot be answered from the released files, which carry no per-record
class-support counts (\texttt{BLOCKERS.md}, ESC-6); we flag it as the follow-up a
dedicated, pre-registered analysis would pursue rather than reporting an untraceable
correlation here.

\subsection{Clinical interpretation}

The method is not an autonomous diagnostic system. It assumes supervised labels
for an early patient segment and evaluates heartbeat classification on curated
research databases. The strong patient heterogeneity---a few records personalize
dramatically, most gain modestly, and a minority regress---is itself a clinical
finding: personalization is unevenly beneficial, and a deployment would need a
per-patient safeguard rather than a blanket ``adaptation always helps''
assumption. The results are algorithmic evidence about memory allocation, not
clinical validation.
